\documentclass[11pt]{article}

\usepackage[margin=1.2in]{geometry}
\usepackage{amsmath,amssymb,amsthm}
\usepackage{enumitem}
\usepackage{booktabs}
\usepackage{xcolor}
\usepackage[colorlinks=true,linkcolor=blue!50!black,citecolor=blue!50!black,urlcolor=blue!50!black]{hyperref}

\newcommand{\F}{\mathbb{F}}
\newcommand{\eq}{\mathsf{eq}}
\newcommand{\RS}{\mathsf{RS}}
\newcommand{\eps}{\varepsilon}
\newcommand{\coef}{\mathrm{coef}}
\newcommand{\What}{\widehat W}

\newtheorem{theorem}{Theorem}
\newtheorem{lemma}[theorem]{Lemma}
\newtheorem{remark}[theorem]{Remark}

\title{Affine Folding over $\F_{2^{128}}$ for Ligerito:\\
Dropping the $\F_{2^{256}}$ Fold Challenges}
\author{DRAFT --- design note}
\date{November 2, 2026}

\begin{document}
\maketitle

\begin{abstract}
Flock reaches 128-bit security in Ligerito's proximity term by sampling
every fold challenge from the quadratic extension $\F_{2^{256}}$ and
splitting each folded extension table into two base-field coordinate
tables at the code switch (\texttt{pcs/ligerito/extension.rs}). The
draft by Chiesa, Yogev, Weissenberg and Zheng, \emph{Reed--Solomon
Proximity Gaps with Small Error via Affine Mixing}, proves mutual
correlated agreement (MCA) for the combination $f + (\alpha + \beta
x)\,g$ with error about $\ell/q$ (the Johnson list size over the field
size). No factor of $n$ appears. This note specifies a Ligerito fold
round that uses affine mixing with all challenges in $\F_{2^{128}}$. It
gives the soundness argument for the three steps the paper leaves
open: claim transport, degree growth and interleaving. It also gives
per-level parameters for the $m = 32$ Fast profile and an
implementation plan. The expected gain is about half of the recursive
commit cost, plus smaller proofs.
\end{abstract}

\section{Setting}\label{sec:setting}

$\F = \F_q$ with $q = 2^{128}$. $D \subseteq \F$ is an
$\F_2$-subspace of size $n$, and $\RS[D, d]$ is the evaluation code of
polynomials of degree at most $d$ on $D$. Messages are coefficient
vectors in the LCH novel basis $(X_j)_{j \ge 0}$, where $\deg X_j = j$.
A vector of length $L$ therefore stores a polynomial of degree less
than $L$, and the code does not depend on the basis.

A Ligerito level commits $2^k$ rows $f(i, \cdot) \in \F^{2^c}$, $i \in
\{0,1\}^k$, as one interleaved codeword. Each Merkle leaf holds the
$2^k$ symbols at one $x \in D$. The level enters with a linear claim
\begin{equation}\label{eq:claim}
  \textstyle\sum_{i, c} f(i, c)\, W(i, c) = v,
\end{equation}
where the weight $W$ is a sum of eq-tensors (the batched opening,
OOD and consistency claims). The level runs $k$ sumcheck rounds over
the row bits. Today every round folds $f' = f_0 + r(f_0 + f_1)$ with
$r \in \F_{2^{256}}$. Here $f_b := f(b, \cdot)$ splits on the bit
bound in that round.

\paragraph{Why $\F_{2^{256}}$ today.} Over $\F_{2^{128}}$, the scalar
line MCA bounds near the Johnson radius carry a factor of at least $n$
(already $n/q$ in the unique-decoding regime). At $n = 2^{20}$ that
leaves about $108$ bits per fold challenge, too many to grind away.
The extension makes the proximity term about $2^{-209}$. The cost is
$\F_{2^{256}}$ arithmetic in every fold and a coordinate split that
doubles every committed recursive table
(\S\ref{sec:cost}).

\section{The affine fold round}\label{sec:round}

Write $Xf_1$ for the table whose rows are the rows of $f_1$ multiplied
by $X$ as polynomials. Each row gains one coefficient. Extend $W_b :=
W(b, \cdot)$ by zero on the new coordinate. One round:

\begin{enumerate}[leftmargin=2em]
\item \textbf{Prover} sends the grid $G[u][b] = \langle u, W_b\rangle$
  for $u \in \{f_0, f_1, Xf_1\}$ and $b \in \{0,1\}$. The claim fixes
  $G[f_1][1] = v - G[f_0][0]$, so the prover sends five
  $\F_{2^{128}}$ elements. Today's round message is two
  $\F_{2^{256}}$ elements, which is four $\F_{2^{128}}$ elements.
\item \textbf{Verifier} samples $(\alpha, \beta, r) \in \F^3$ from one
  fused squeeze, with the fold grinding of \S\ref{sec:params}.
\item \textbf{Both} move to
  \[
    f' = f_0 + \alpha f_1 + \beta\, Xf_1, \qquad
    W' = W_0 + r(W_0 + W_1),
  \]
  \[
    v' = \textstyle\sum_{u, b} a_u(\alpha, \beta)\, e_b(r)\, G[u][b],
    \qquad (a_{f_0}, a_{f_1}, a_{Xf_1}) = (1, \alpha, \beta), \quad
    (e_0, e_1) = (1 + r, r).
  \]
\end{enumerate}
By bilinearity, $\langle f', W'\rangle = v'$ when the grid is honest.
The weight update is today's binding at $r$, so $W'$ keeps its
eq-structure and the verifier evaluates it as it does now. The
verifier never evaluates an $X$-shifted weight.

\paragraph{Codeword view.} For every $x \in D$,
\[
  \mathrm{Enc}(f')(x) = \mathrm{Enc}(f_0)(x) + (\alpha + \beta x)\,
  \mathrm{Enc}(f_1)(x),
\]
which is the paper's affine combination applied row by row. After the
$k$ rounds of a level, a query at $x$ combines the opened leaf with the
weights $\omega_i(x) = \prod_{j=1}^{k} (\alpha_j + \beta_j x)^{i_j}$.
These weights depend on $x$, so the verifier computes $2^k$ products
per query plus the $2^k$-term inner product. Today that inner product
is $\F_{2^{128}} \times \F_{2^{256}}$, which costs two base
multiplications per term. The count of base multiplications is
therefore unchanged.

\paragraph{The multiply-by-$X$ map.} Let $W_i$ be the subspace
vanishing polynomials of the domain basis $(v_i)$ and $\What_i = W_i /
W_i(v_i)$. The identity $W_{i+1}(X) = W_i(X)\,(W_i(X) + W_i(v_i))$
gives
\[
  \What_i^{\,2} = a_i\, \What_{i+1} + \What_i, \qquad a_i =
  W_{i+1}(v_{i+1}) / W_i(v_i)^2,
\]
and $X = v_0 \What_0$. With the standard basis $v_i = x^i$ we have
$v_0 = 1$ and $X = \What_0$. Write $P_i = a_0 \cdots a_{i-1}$, and let
$t(j)$ and $z(j)$ be the numbers of trailing ones and trailing zeros of
$j$. Expanding $\What_0 X_j$ through the carries gives
$X \cdot X_j = \sum_{i < t(j)} P_i\, X_{j - 2^i + 1} + P_{t(j)}\,
X_{j+1}$. In gather form, each output coefficient has exactly two
terms:
\[
  (T g)_m = P_{t(m-1)}\, g_{m-1} + P_{z(m)}\, g_{m + 2^{z(m)} - 1}
  \qquad (m \ge 1;\ (Tg)_0 = 0),
\]
where the second term is present when its index is in range. The
prover applies $T$ in $O(n)$ with two multiplications by constants per
coefficient. \texttt{benches/affine\_fold\_probe.rs} checks this
formula against the NTT: $\mathrm{Enc}(Tg)(x) = x\,\mathrm{Enc}(g)(x)$
at every domain point. The prover
needs $Xf_1$ for $f'$ and for the two grid entries $\langle Xf_1,
W_b\rangle$.

\section{Soundness}\label{sec:soundness}

Each fold round contributes two errors to the round-by-round
analysis: the MCA error of the affine combination
(\S\ref{sec:mca}) and the claim-transport error of the grid
(\S\ref{sec:transport}).

\subsection{Claim transport}\label{sec:transport}

\begin{lemma}[Grid transport]\label{lem:transport}
Fix tables $f_0, f_1$ and weights $W_0, W_1$. Suppose the incoming
claim is false: $v \ne \langle f_0, W_0\rangle + \langle f_1,
W_1\rangle$. Then for any grid the prover sends,
$\Pr_{\alpha, \beta, r}[\, v' = \langle f', W'\rangle \,] \le 2/q$.
\end{lemma}

\begin{proof}
Let $\Delta[u][b]$ be the sent grid minus the true inner products. The
claim relation forces $\Delta[f_0][0] + \Delta[f_1][1] \ne 0$, so
$\Delta \ne 0$. The difference $v' - \langle f', W'\rangle =
\sum_{u, b} a_u e_b \Delta[u][b]$ is a polynomial in $(\alpha, \beta,
r)$ of total degree at most $2$. Its six products $\{1, \alpha,
\beta\} \times \{1 + r, r\}$ are linearly independent, so a nonzero
$\Delta$ gives a nonzero polynomial. Schwartz--Zippel gives at most
$2/q$.
\end{proof}

When the analysis carries $L$ candidate codewords through the round,
the transport error is $2L/q$. Under the Johnson+OOD regime that Flock
uses, OOD reduces the list to one candidate, so $L = 1$. This term
replaces today's sumcheck term (about $250$ bits over
$\F_{2^{256}}$). It is about $127$ bits here and is absorbed into the
fold grinding.

\subsection{MCA for one round}\label{sec:mca}

\paragraph{Degree growth.} Round $j$ of a level combines words of
$\RS[D, d + j - 1]$ into $\RS[D, d + j]$ on the same domain, with
$d = 2^c - 1$. Both halves have degree at most $d + j - 1$, so the
paper's condition $\deg G \le \deg_{\mathrm{out}} - 1$ holds with no
change. The intermediate codes appear only in the analysis. Nothing is
committed at those degrees, and the rate grows by $k/n$, which is
negligible. The next level re-encodes the folded message at its own
rate (\S\ref{sec:overflow}), so the degree does not accumulate across
levels.

\paragraph{Interleaving (Habo\v{c}k's observation).} Round $j$
combines $t = 2^{k-j}$ row pairs with one shared $(\alpha, \beta)$,
and the rows must share an agreement set. Because $D \subseteq \F_q$,
an interleaved word of $t$ rows of $\RS_q[D, d]$ is exactly one word of
$\RS_{q^t}[D, d]$: fix an $\F_q$-basis of $\F_{q^t}$. The rate and
distance are the same, and agreement of the $\F_{q^t}$-word is
simultaneous agreement of all rows. Since $x \in D \subseteq \F_q$ and
$(\alpha, \beta) \in \F_q^2$, the affine combination of
$\F_{q^t}$-words is the row-wise combination.

\begin{lemma}[Subfield challenges]\label{lem:subfield}
Let $\F_q \subseteq \F_Q$, $D \subseteq \F_q$ and $f, g: D \to \F_Q$.
The bound of the paper's Lemma~3.1 on the number of exceptional
$(\alpha, \beta)$ holds for $(\alpha, \beta) \in \F_q^2$, with $q$
replaced by $Q$ only inside the $\max\{q, \binom n2\}$ factor, where it
stays $q$. Hence
\[
  \Pr_{(\alpha, \beta) \leftarrow \F_q^2}[\text{exceptional}] \le
  \frac{(4d^2M^2 + 1)D^3 v^2 + \ell \max\{q, \binom n2\}}{q^2}.
\]
\end{lemma}

\begin{proof}[Proof sketch]
The paper's proof counts bad pairs and divides by $q^2$ at the end.
We check each count with words over $\F_Q$ and pairs from $\F_q^2$.
\begin{itemize}[leftmargin=1.5em]
\item \emph{Lemma 2.3, rational case.} The good pairs form a proper
  affine subspace of $\F_Q^2$. It meets $\F_q^2$ in at most $q$
  points. In the polynomial case, two coordinates still fix at most
  one pair.
\item \emph{Lemma 3.3.} The excluded pairs give distinct linear factors
  $Y - \alpha' - \beta' X$ of a fixed polynomial, with $\alpha',
  \beta' \in \F_q \subseteq \F_Q$. The Frobenius root map is a
  bijection on $\F_q$. The count $D^2$ is unchanged.
\item \emph{Lemma 3.4.} The list size $\ell = L_{d+1}(B-1)$ is a
  set-system (Johnson) bound and does not depend on the field.
  Directions $(\sigma, \omega, H)$ range over $\F_Q$, which only
  enlarges the persistent case that the lemma already covers.
\item \emph{Claims 3.6 and 3.7.} Take $K$ to be the algebraic closure
  of $\F_Q(T_0)$. The incidence and degree bounds are geometric and do
  not depend on the field. Hensel lifting needs $R_Z(X, \alpha' +
  \beta' X, P(X^e)) \ne 0$ over $\F_Q$, and $T_0$ is transcendental
  over $\F_Q$.
\end{itemize}
\end{proof}

\begin{remark}
Lemma~\ref{lem:subfield} needs to be checked by the authors. Ulrich
Habo\v{c}k has offered a write-up of the interleaved argument for
lines. Applying the lemma round by round to $t = 2^{k-1}, \dots, 1$
pairs gives the multilinear fold with error $\eps_{\mathrm{aff}}$ per
round and \emph{no} union over rows. Today's analysis pays $2^{k-1}$
there (the row union in the mca-commutes lemma).
\end{remark}

\paragraph{Agreement loss.} The theorem allows agreement loss $s = 1$
per round, so up to $k$ coordinates per level. The query bound
$a^t$ becomes $(a + k/n)^t$, a factor of at most $\exp(tk/(an))$. At
L0 ($t = 244$, $k = 6$, $n = 2^{20}$, $a \approx 0.727$) this is
$0.003$ bits. At L5 ($t = 25$, $k = 4$, $n = 2^{15}$, $a \approx
0.042$) it is $0.10$ bits. The query budget must include it.

\subsection{Overflow coefficients}\label{sec:overflow}

After $k$ rounds the folded message has degree at most $d + k$. Its
coefficients on $X_{2^c}, \dots, X_{2^c + k - 1}$ are $\tau_1, \dots,
\tau_k$. At the code switch the prover sends the $\tau_s$ in the clear,
and the next level commits only the low $2^c$ coefficients. The claim
weight is zero on the overflow coordinates, so the claim
$\langle f', W'\rangle$ does not change. Each consistency claim at a
query $x$ becomes
\[
  \textstyle\langle f'_{\mathrm{lo}}, (X_c(x))_c\rangle =
  \sum_i \omega_i(x)\, \mathrm{leaf}_i(x) - \sum_s \tau_s\,
  X_{2^c + s - 1}(x).
\]
The verifier evaluates $X_{2^c + s - 1}(x)$ from products of
$\What_i(x)$, which are $\F_2$-linear in $x$.

\section{Parameters}\label{sec:params}

Per fold round, $\eps_{\mathrm{round}} \le \eps_{\mathrm{aff}} + 2/q$,
and the fold grinding is $\lceil 128 - \log_2
(1/\eps_{\mathrm{round}}) \rceil$ bits. The script
\texttt{scripts/affine\_fold\_params.py} evaluates the general form of
the paper's Lemma~3.1 for every round of every level of a config. It
uses degree $d_j = d + j$ in round $j$, the integer Johnson bound for
$\ell$, and a search over $m$, $D$ and $B$ subject to the paper's
inequality (3.4). Any $(m, D, B)$ it reports is a valid certificate,
so each bound is proven. By Lemma~\ref{lem:subfield} the bound does not
depend on the number of interleaved rows. The closed form of the
paper's Theorem~1.1 does not apply at levels L2--L5, because
$n\eta^2 < 4/\rho$ there. Lemma~3.1 still does.

The table shows the output for \texttt{configs/ligerito/m32\_fast.toml}
($\eta = 0.02$), taking the worst round of each level:

\begin{center}
\begin{tabular}{@{}lcccccc@{}}
\toprule
Level & rate $\rho$ & $n$ & $\ell$ & $\log_2 \eps_{\mathrm{aff}}$ &
fold grind (bits) & agreement loss (bits) \\
\midrule
L0 & $2^{-1}$  & $2^{20}$ & 7     & $-125.2$ & 4 & 0.003 \\
L1 & $2^{-3}$  & $2^{19}$ & 16    & $-124.0$ & 5 & 0.002 \\
L2 & $2^{-5}$  & $2^{18}$ & 21    & $-123.6$ & 5 & 0.005 \\
L3 & $2^{-7}$  & $2^{17}$ & 22    & $-123.5$ & 5 & 0.014 \\
L4 & $2^{-9}$  & $2^{16}$ & 24    & $-123.4$ & 5 & 0.040 \\
L5 & $2^{-11}$ & $2^{15}$ & 28--30 & $-123.1$ & 5 & 0.104 \\
\bottomrule
\end{tabular}
\end{center}

Grinding is fused into the challenge squeeze, so $4$--$5$ bits per fold
challenge cost little. A production analysis should emit a new
\texttt{analysis\_version} from these numbers.

\section{Expected cost}\label{sec:cost}

Ligerito prover sub-phases at $m = 32$ (Fast profile, main
\texttt{2c735c3}, Apple M4 Max under background load, one warm run):

\begin{center}
\begin{tabular}{@{}lrl@{}}
\toprule
Sub-phase & ms & Effect of the affine fold \\
\midrule
recursive commits (NTT + Merkle) & 28.5 & about half: no coordinate split \\
\texttt{initial\_k} folds + code switch & 12.7 & about flat; the split pass goes away \\
recursive folds + code switches & 4.2 & smaller ($\F_{2^{128}}$ vs $\F_{2^{256}}$) \\
induce (basis + sum) & 4.7 & smaller ($\F_{2^{128}}$ basis) \\
OODs, grinding, opens & 5.9 & flat \\
\midrule
total & 55.0 & estimate $-15$ to $-20$ ms of a $248$ ms prove \\
\bottomrule
\end{tabular}
\end{center}

\paragraph{Probe measurement.} \texttt{benches/affine\_fold\_probe.rs}
times the real \texttt{ligero\_commit} at today's split shapes and at
the half-size native shapes, keeping the same rates and four lanes. Run
on the same machine (median of 7, ms):

\begin{center}
\begin{tabular}{@{}lrrrr@{}}
\toprule
Level & split table & ms & native table & ms \\
\midrule
L1 & $2^{20}$ & 15.47 & $2^{19}$ & 8.58 \\
L2 & $2^{17}$ & 7.43 & $2^{15}$ & 2.08 \\
L3 & $2^{14}$ & 3.72 & $2^{11}$ & 0.80 \\
L4 & $2^{11}$ & 2.01 & $2^{7}$ & 0.32 \\
L5 & $2^{8}$ & 1.06 & --- & --- \\
\midrule
total & & 29.70 & & 11.78 \\
\bottomrule
\end{tabular}
\end{center}

The commits alone save about $18$ ms, close to the trace's
$28.5$ ms total for today's recursive commits. The probe also runs
simple fold kernels with a dense basis, written the same naive way for
both paths. At L0 ($2^{25}$, $k = 6$) the F256 fold takes $89.6$ ms and
the affine fold $85.0$ ms. At L1 the F256 fold on $2^{20}$ takes
$3.9$ ms and the affine fold on $2^{19}$ takes $2.2$ ms. Production's
fused F256 kernels run the L0 folds in $12.7$ ms, so the affine path
needs the same fusion work (step~3 of \S\ref{sec:plan}) to keep that
part flat.

\paragraph{Prototype.} \texttt{pcs/ligerito/affine.rs} implements the
full affine ladder on the plain native path: the grid, the $T$ map,
the affine fold, overflow coefficients, OOD and consistency glue,
native commits, and the verifier. Its tests check $T$ against the NTT
and run a roundtrip and tamper cases. \texttt{benches/affine\_vs\_f256.rs}
gives both provers the same $2^{25}$ witness, dense eq basis and L0
commitment ($2^6$ lanes, rate $1/2$), in interleaved runs. The F256
side keeps its fused and lookahead kernels; the affine side has only
the simple ones.
\begin{center}
\begin{tabular}{lrrr}
\toprule
& F256 split & affine & change \\
\midrule
prove, run 1 (ms, median of 9) & 107.7 & 83.3 & $-22.7\%$ \\
prove, run 2 (ms, median of 9) & 103.4 & 81.7 & $-21.0\%$ \\
proof size (KiB) & 493.8 & 465.3 & $-5.8\%$ \\
recursive levels & 5 & 4 & \\
\bottomrule
\end{tabular}
\end{center}
In one warm trace the recursive commits drop from $31.3$ to $20.9$ ms,
and the folds take $43.3$ ms against $45.8$ ms for the fused F256
folds. The affine verifier runs in $13.4$ ms. This is the native path,
not the production open, which uses a lane-major 46-lane L0 and a
virtual basis.

\paragraph{The production L0.} The production open commits $46$ live
lanes of $2^{19}$ columns, lane-major, and its weight is a sum of eq
tensors that is never materialized. The F256 ladder reads the witness
once for per-lane statistics $\langle \mathrm{lane}_i, e\rangle$, where
$e$ is the column factor of a term. All six L0 rounds then run on
these $64$-entry vectors. The seeded term's statistics come from the
merged sumcheck, so only the OOD term is swept. The affine round also
needs $\langle X f_1, W_b\rangle$, which is a dot product with
$(T^{\top})^s e$ for $s = 1, \dots, 6$. As plain dot products that is
$13$ sweeps instead of one. They measured $25$ ms against $2.3$ ms. A
NEON kernel with three PMULLs per product did no better, because
the sweep is PMULL-bound.

One fold of the row gives all $s$ at once. For $g$ with $2^c$
coefficients and $z \in \F^c$, let $F_s(g) = \langle (X^s
g)_{<2^c}, \mathrm{eq}(z)\rangle$. Write $g = g_L + \What_{c-1}\, g_H$.
Multiplication by $X$ raises a coefficient index by at most one, and
coefficients at $2^c$ and above never move back below $2^c$. With
$\What_{c-1}^2 = a_{c-1}\What_c + \What_{c-1}$ this gives, for $s \le
2^{c-1}$,
\[
  F_s(g) = F'_s\bigl(g_L + z_{c-1}(g_L + g_H)\bigr)
  + z_{c-1}\, \langle B_s(g_L + g_H), \mathrm{eq}(z_{<c-1})\rangle,
\]
where $F'_s$ is the same functional one level down and $B_s(h)$ is the
overflow of $X^s h$ past $2^{c-1}$. $B_s(h)$ depends only on the top
$s$ coefficients of $h$. The fold costs what $F_0$ alone costs, and
the corrections cost $O(c\, s^2)$ per lane. The prover folds the top
$7$ column bits in one pass against their eq table, gathers the few
boundary coefficients for those levels, and recurses on the remaining
$2^{12}$ words. The rounds then update the statistics in place:
$S'[m][s] = S[2m][s] + \alpha S[2m+1][s] + \beta S[2m+1][s+1]$. After
L0 the prover materializes $\sum_i \omega_i(X)\,\mathrm{lane}_i$ in
one column-chunked pass and a Horner loop over $T$. Tests check the
identity against brute force and the fused sweep against the per-lane
recursion.

\texttt{benches/affine\_vs\_f256.rs lanes} runs both provers on this
shape: \texttt{commit\_lane\_major} with $46$ lanes, the factored eq
basis, and the seeded statistics. The F256 side is the production
statistics ladder. Two runs of nine interleaved pairs, under some
background load:
\begin{center}
\begin{tabular}{lrrr}
\toprule
& F256 ladder & affine & change \\
\midrule
prove, run 1 (ms, median) & 58.2 & 40.5 & $-30.4\%$ \\
prove, run 2 (ms, median) & 60.5 & 39.3 & $-35.0\%$ \\
prove, best of 9 & 56.3 & 38.3 & $-32.1\%$ \\
proof size (KiB) & 425.2 & 396.6 & $-6.7\%$ \\
verify (ms) & & 1.4 & \\
\bottomrule
\end{tabular}
\end{center}
In a warm trace the affine L0 takes $16.6$ ms ($5.5$ statistics,
$11.1$ materialize) against $10.9$ ms for the F256 ladder. The
recursive commits take $11.9$ ms against $29.5$ ms, and the
consistency step $3.9$ ms against $5.7$ ms. That step had cost
$10.3$ ms before the subspace constants $s_k(v_k)^{-1}$ were cached
(each query inverted about $20$ of them) and the queries ran in
parallel. The same fix took the verifier from $13.5$ to $1.4$ ms.

The materialize pass does $\sum_i (\mathrm{popcount}(i) + 1) = 165$
multiplications per column against $92$ for the F256 fold, at about
the same rate per product, so it stays the largest L0 cost. Two
attempts did not help. Folding the low lane bits pairwise inside the
column chunks cuts the count to about $150$, because $T$ reads outside
an aligned chunk only at its first column, but it ran no faster. A
register-resident Karatsuba accumulator with pre-split coefficients ran
slower than the plain loop.

The largest item is the split. Today a folded $\F_{2^{256}}$ table of
length $N$ is committed as a $\F_{2^{128}}$ table of length $2N$, and
each level removes only $k - 1$ variables. With the affine fold, a
table of length $N$ is committed as length $N$ (plus $k$ overflow
values), and each level removes $k$ variables. The ladder should be
retuned. Proofs also shrink: round messages, OOD answers and the final
message become $\F_{2^{128}}$, and the recursive Merkle leaves halve.
The 100-bit profiles benefit as well, since they pay for
$\F_{2^{256}}$ today with no security gain.

\section{Implementation plan}\label{sec:plan}

\begin{enumerate}[leftmargin=2em]
\item \textbf{Analysis.} Commit the parameter script and the analysis
  identity. Get Lemma~\ref{lem:subfield} checked by the paper's
  authors and by Habo\v{c}k.
\item \textbf{Prototype.} Add \texttt{pcs/ligerito/affine.rs} next to
  \texttt{extension.rs}, behind a new profile field value, for the
  native opening only. Start with simple kernels: the grid message,
  the $T$ map, the affine fold, overflow handling and per-query
  weights. Measure open time and proof size at $m = 32$ against
  \texttt{extension.rs} with interleaved runs.
\item \textbf{Kernels.} If the gain holds, port the fast paths. The L0
  lookahead and skip kernels and the fused fold+message passes assume
  multilinear folding, so they need affine versions. This step is
  where most of the engineering time goes. \emph{Status:} the
  lane-major L0 with the factored basis is done (statistics-only
  rounds, \S\ref{sec:cost}). Open: the merged-transport wiring
  (\texttt{pcs::open\_batch\_merged}), and the recursive-level
  kernels.
\item \textbf{Recursion and GPU.} Port the in-circuit verifier
  (\texttt{tower/query.rs}, \texttt{gates\_spine.rs}, the walkers):
  remove the $\F_{2^{256}}$ gadgets and add the grid check, per-query
  weights and overflow terms. Port the CUDA twin, re-pin fixtures,
  bump the proof-IO version and retire \texttt{extension.rs}.
\end{enumerate}

\paragraph{Open items.}
(i) Lemma~\ref{lem:subfield} as a checked statement.
(ii) Round-by-round accounting of the agreement loss together with the
OOD list reduction.
(iii) The overflow handling at L0, where the leaf is lane-major and has
46 lanes.
(iv) Whether the per-query weights $\omega_i(x)$ change the
recursion circuit's row count enough to cross a power of two.

\paragraph{References.}
A.~Chiesa, E.~Yogev, G.~Weissenberg, K.~Z.~Zheng, \emph{Reed--Solomon
Proximity Gaps with Small Error via Affine Mixing}, draft, 2026.
B.~B\"unz, R.~Rothblum, W.~Wang, \emph{Flock}, ePrint 2026/1329,
Appendix~C (Ligerito; mca-commutes lemma).
U.~Habo\v{c}k, private communication on interleaved correlated
agreement, 2026.

\end{document}
